\section{Experimental Setup}

\subsection{Dataset}

\textbf{Symbol}: SPY (S\&P 500 ETF)

\textbf{Primary Dataset (2024)}:
\begin{itemize}
\item Period: January 2, 2024 - December 31, 2024
\item Trading days: 252
\item 30-day windows: 223 (rolling windows with 29-day stride)
\item Data source: Alpha Vantage Premium API (options chain + spot prices)
\end{itemize}

\textbf{Comparison Dataset (2020)}:
\begin{itemize}
\item Period: January 2, 2020 - December 31, 2020
\item Trading days: 252
\item 30-day windows: 223
\item Purpose: Pre-0DTE baseline for regime persistence comparison
\end{itemize}

\textbf{Negative Control Datasets}:
\begin{itemize}
\item Shuffle test: 404 windows (2024 Q1 + 2020, randomized day order)
\item Transitional test: 202 windows (2024 Q1 + 2020, high flip counts)
\item Low-magnitude test: 203 windows (2024 Q1 + 2020, scaled <\$5B)
\end{itemize}

All options data includes strike price, open interest, implied volatility, days to expiration, and contract type (call/put). Spot prices obtained from daily close values.

\subsection{Data Processing Pipeline}

\subsubsection{GEX Calculation}

For each trading day:
\begin{enumerate}
\item Fetch options chain (all strikes, all expirations)
\item Calculate Black-Scholes gamma using implied volatility from market prices
\item Weight by open interest: $\text{GEX}_{\text{strike}} = \text{OI} \times \Gamma \times S^2 \times 0.01 \times 100$
\item Aggregate across strikes: $\text{GEX}_{\text{total}} = \sum_{\text{calls}} \text{GEX} - \sum_{\text{puts}} \text{GEX}$
\item Store net GEX value (positive = dealer long gamma, negative = dealer short gamma)
\end{enumerate}

\subsubsection{Window Generation}

Generate 30-day rolling windows with stride=1:
\begin{itemize}
\item Window 1: Days 1-30
\item Window 2: Days 2-31
\item Window 3: Days 3-32
\item \ldots
\item Window N: Days (N)-(N+29)
\end{itemize}

For each window, calculate:
\begin{itemize}
\item Persistence: $\frac{\max(\text{positive days}, \text{negative days})}{30}$
\item Magnitude: $\frac{1}{30}\sum_{i=1}^{30} |\text{GEX}_i|$
\item Flips: $\sum_{i=1}^{29} \mathbb{1}[\text{sign}(\text{GEX}_{i+1}) \neq \text{sign}(\text{GEX}_i)]$
\end{itemize}

\subsubsection{Obfuscation}

Transform each window before LLM submission:
\begin{itemize}
\item Replace calendar dates with relative labels: "Day T-29", "Day T-28", \ldots, "Day T+0"
\item Replace "SPY" with generic label: "INDEX\_1"
\item Remove all temporal context (day of week, month, quarter, year)
\item Present only: GEX sequence (30 values), spot price, current date label
\end{itemize}

\subsection{Sample Size and Window Breakdown}

The five-phase validation structure tests overlapping windows across multiple analyses. Table~\ref{tab:windows_breakdown} reconciles the total sample sizes:

\begin{table}[h]
\centering
\caption{Comprehensive Window Sample Breakdown Across All Validation Phases}
\small
\begin{tabular}{lrrrl}
\toprule
\textbf{Phase} & \textbf{Windows} & \textbf{Type} & \textbf{Overlap w/ 4A} & \textbf{Unique} \\
\midrule
Phase 1 (Q1 2024) & 52 & Real & ⊂ Phase 3 & 0 \\
Phase 2a (Shuffle) & 404 & Synthetic & No & 404 \\
Phase 2b (Transitional) & 202 & Synthetic & No & 202 \\
Phase 2c (Low-Magnitude) & 203 & Synthetic & No & 203 \\
Phase 3 (Full 2024) & 223 & Real & ⊂ Phase 5 & 0 \\
Phase 4 (Full 2020) & 223 & Real & ⊂ Phase 5 & 0 \\
Phase 5 (2020-2025) & 1,412 & Real & -- & 1,412 \\
\midrule
\textbf{Total Real Windows (2020-2025)} & -- & -- & -- & 1,412 \\
\textbf{Total Synthetic Controls} & -- & -- & -- & 809 \\
\textbf{Total Window Evaluations} & -- & -- & -- & 2,221 \\
\bottomrule
\end{tabular}
\label{tab:windows_breakdown}
\end{table}

\noindent\textbf{Explanation}: The abstract references ``five phases across six years (2020-2025, 1,858 windows)'' referring to the rolling 30-day windows available from historical trading data. The 1,858 figure represents the sum of unique real windows: Phase 5 contains all 2020-2025 windows (1,412), with Phases 1, 3, and 4 being subsets. When considering only unique real windows, the total is 1,412 from Phase 5. However, the five phases test these windows multiple times for robustness. Phase 1 tests a subset to establish baseline detection. Phase 3 tests the full 2024 year. Phase 4 tests the full 2020 year. Phase 5 re-tests all 1,412 unique windows (2020-2025) as part of the multi-year analysis.

Synthetic negative controls (Phase 2: 809 windows) are generated from the real data by shuffling, filtering for transitional regimes, or scaling for low magnitude. These 809 synthetic windows are additional to the 1,412 real windows, yielding 2,221 total evaluations.

The five-phase structure prioritizes robustness through repeated testing and negative controls, even though the unique real-world data represents 1,412 trading windows.

\subsection{LLM Configuration}

\textbf{Model}: OpenAI o4-mini-2025-04-16 (reasoning model)

\textbf{API Parameters}:
\begin{itemize}
\item Temperature: 1.0 (default reasoning temperature)
\item Max tokens: 16,384 (sufficient for detailed chain-of-thought)
\item Top-p: 1.0 (no nucleus sampling)
\end{itemize}

\textbf{Batch Processing}:
\begin{itemize}
\item Method: OpenAI Batch API (asynchronous processing)
% \item Cost savings: 50\% vs synchronous API (\$0.075 vs \$0.15 per 1M input tokens) % Omitted - not tracking cost
\item Processing time: 6-23 minutes per batch (223 windows)
\item Reliability: 100\% completion rate (no timeouts or rate limits)
\end{itemize}

\subsection{Prompt Structure}

Each prompt contains:

\vspace{0.2cm}

\noindent\textbf{System Message}:
\begin{verbatim}
You are a financial market analyst identifying
persistent dealer gamma regimes from 30-day
GEX sequences.
\end{verbatim}

\vspace{0.2cm}

\noindent\textbf{User Message}:
\begin{itemize}
\item 30-day GEX sequence with obfuscated dates
\item Regime classification criteria (70\% persistence, \$5B magnitude, $\leq$5 flips)
\item Explicit instructions for step-by-step reasoning
\item Output format: JSON with regime\_type, confidence, reasoning, metrics
\end{itemize}

\vspace{0.2cm}

\noindent\textbf{Output Schema}:
\begin{verbatim}
{
  "regime_type": "persistent_negative",
  "regime_detected": true,
  "confidence": 85,
  "reasoning": "...",
  "llm_metrics": {
    "persistence_pct": 96.67,
    "avg_magnitude_b": 13.45,
    "sign_flips": 1
  }
}
\end{verbatim}

\subsection{Validation Protocol}

Our multi-phase validation strategy tests both framework selectivity and market structure evolution:

\vspace{0.2cm}

\noindent\textbf{Phase 1 - Q1 2024 Baseline} (52 windows): Establish baseline detection rate on January-March 2024 data. Success criteria: 30-50\% detection demonstrating selectivity rather than universal pattern matching.

\vspace{0.2cm}

\noindent\textbf{Phase 2 - Negative Controls} (809 windows): Validate discrimination through three synthetic tests. Shuffle test (404 windows) randomizes day order within sequences to destroy temporal structure, expected $<$20\% false positive rate. Transitional test (202 windows) filters high-volatility periods with 7-10 sign flips, expected $<$10\% FP. Low-magnitude test (203 windows) scales persistent sequences below \$5B threshold, expected $<$10\% FP.

\vspace{0.2cm}

\noindent\textbf{Phase 3 - Full 2024 Validation} (223 windows): Test generalization across all 252 trading days in 2024 to determine whether Q1 results extend to full year or reflect quarterly anomaly.

\vspace{0.2cm}

\noindent\textbf{Phase 4 - 2020 Comparison} (223 windows): Test pre-0DTE era (2020, 0DTE $<$5\% volume) versus post-0DTE era (2024, 0DTE $\approx$48\% volume). Hypothesis: Lower detection rate in 2020 indicates 0DTE increased regime persistence.

\subsection{Quality Control}

\vspace{0.2cm}

\noindent\textbf{Mechanical Verification}:

For each LLM response, we extract the model's stated metrics (persistence\_pct, avg\_magnitude\_b, sign\_flips) and compare against ground truth calculated from actual GEX data. Windows with >5\% metric discrepancy are flagged for manual review.

\vspace{0.2cm}

\noindent\textbf{Output Validation}:

All LLM outputs validated for correct structure and required fields:
\begin{itemize}
\item Valid JSON structure (LLM responses)
\item Required fields present (regime\_type, confidence, reasoning, llm\_metrics)
\item Confidence in valid range [0, 100]
\item Regime type in valid set (persistent\_positive, persistent\_negative, transitional, low\_conviction)
\end{itemize}

Validation results stored in YAML format for reproducibility and batch metadata in JSON. Invalid responses logged and excluded from analysis (Phase 3: 0\% invalid, Phase 4: 0\% invalid).

\subsection{Reproducibility}

All code, data preprocessing scripts, and validation results are available in the public GitHub repository: \url{https://github.com/iAmGiG/gex-llm-patterns}. Validation runs can be reproduced using the exact prompt templates and experimental configurations provided in the repository.
